\begin{abstract}
本稿は、ブロックストレージ装置内で観測可能な安価な I/O 統計から
ランサムウェアらしい異常を早期検出できるかを検証するための
研究プロトコルである。現段階の貢献は性能値ではなく、観測境界、
異常分類、入力スキーマ、ベースライン、アブレーション、証拠ゲートを
査読可能な形で固定することにある。対象装置は SCSI/NVMe 的な
read/write コマンド、LBA またはオフセット、転送長、時刻、および
既存の装置カウンタがある場合だけ低コストな圧縮率または
エントロピー相当信号を観測する。ペイロード走査や per-block Shannon
entropy のような重い特徴量は、組み込み主張の前提にしない。

本稿が教師なし異常検出を採用する理由は、既知 family の識別ではなく、
未知または変種のランサムウェアに対して、正常時系列からの逸脱として検出する
余地を残すためである。その中で AE は、固定長 tensor を再構成するだけの
実装に落としやすく、ブロックストレージ装置内の MNN CPU 推論候補として扱いやすい。
中心仮説は、正常な 10 秒 I/O 統計列を学習した小型 AutoEncoder (AE) の
再構成誤差が、write-heavy burst、write ratio の変化、平均 LBA の移動、
平均 I/O 長の変化、圧縮率またはエントロピー相当信号の変化を伴う
攻撃窓で上昇し、単純な write/entropy ルールを超える情報を持つかもしれない、
というものである。AE は既定解ではなく、未知脅威対応と実装容易性を両立し得る
メモリ制約付き仮説として扱う。
入力は 10 秒フレームを $N$ 個束ねた $N \times 10$ 秒の時系列であり、
最終実装では Alibaba MNN による CPU 推論を想定する。ストレージ装置では
多数の volume/namespace を定義できるため、単一 volume では軽い検出器でも
全 volume へ常時適用するとメモリと CPU 時間が累積する。したがってモデル経路に
使えるメモリは volume あたり MNN ランタイムを除いて 500 KB とし、性能、MNN 変換、
score parity、per-volume peak model-owned memory、aggregate memory/CPU は、
manifest 付き実験と測定が
揃うまで非主張とする。
\end{abstract}
